\chapter{Experiments: status and what is settled}
\label{ch:experiments}

This chapter is the current state of the empirical work. It is deliberately explicit about
what is finished, what is running, and what has not been started, because several results here
are recent and one of them changed the framing of the accompanying note.

\section{Settled}

\begin{center}\small
\begin{tabular}{@{}llp{7.4cm}@{}}
\toprule
question & where & finding \\
\midrule
Is the recursion right? & \codefile{exp\_01}, tests &
Agrees with brute-force enumeration to $1.6\times10^{-14}$ and with the closed-form Gaussian
score to $9.2\times10^{-15}$. \\
\addlinespace
What does the Gaussian closure cost? & \codefile{exp\_02}, \codefile{exp\_03} &
Relative score deviation $0.198$ at $t=0.08$ falling to $9.9\times10^{-5}$ at $t=2.4$;
largest at \emph{low} noise. \\
\addlinespace
Can an unknown innovation law be recovered from noised data? & \codefile{exp\_06} &
Yes. A mixture kernel that has never seen a Laplace density recovers its variance essentially
exactly and its excess kurtosis monotonically in $C$, reaching $2.71$ against a true $3.0$ at
$C=8$. \\
\addlinespace
Gradient ascent or EM? & \codefile{exp\_08} &
Neither dominates; see \S\ref{ch:estimation}. \\
\addlinespace
Denoising accuracy against networks & \codefile{exp\_07}, \codefile{replicates} &
Ratio $9$--$14$ across seven doublings of the data, six seeds. \\
\addlinespace
Does a locality-respecting architecture close the gap? & \codefile{exp\_12} &
Most of it. Margin falls from ${\sim}10$ to $2.3$--$4.2$; between $58\%$ and $73\%$ of the
fully connected network's deficit was its architecture. \\
\addlinespace
How deep must a network be to run BP on a chain? & \codefile{exp\_17} &
A chain-local network implements one or two message steps to within $1$--$6\%$ of its exact
information horizon, then plateaus; absolute error bottoms at depth $6$ and rises after. Depth
is not the representational constraint. \\
\bottomrule
\end{tabular}
\end{center}

\section{The result that changed the framing}
\label{sec:dissociation}

\codefile{exp\_16} compares the same estimators \emph{through reverse diffusion} rather than
pointwise. The validation passes first: run with the true kernel, the sampler reproduces the
closed-form target for $P_{t_{\min}}$ to within $1.4\%$ across budgets.

At $\nseq=2048$ on a Laplace chain, the pointwise and distributional rankings \emph{disagree}.
The estimator is an order of magnitude better at denoising and three times better on the
covariance profile, while the convolutional network reproduces the innovation law more
faithfully by both kurtosis and divergence. The fully connected network shows it most sharply:
over the same range its denoising error improves twelvefold while its generated kurtosis
degrades from $0.850$ to $0.206$.

That is the observation. It raises two questions --- whether it is an artefact of this
particular setting, and what causes it --- and \S\ref{sec:resolved} answers both. The reader
who wants the conclusion rather than the sequence can skip there directly; what follows is
worth reading only for how the two candidate explanations were separated.

\section{Resolved since: the dissociation is a capacity effect}
\label{sec:resolved}

Four sweeps have since returned, and together they settle what the dissociation is.

\paragraph{It is not an artefact of the setting.} Across the four non-Gaussian families of
\S\ref{sec:matched-covariance} at matched covariance, the convolutional network is closer to
the target innovation kurtosis than the estimator at $C=4$ in every case. It persists at
$\nsites\in\{32,64,128\}$. And it is \emph{identical to three decimals} under grid refinement
$\ngrid\in\{401,801,1601\}$ --- the strongest of the three controls, since it removes
quadrature as an explanation outright.

\paragraph{It is a limit of the model class, not of the data.} Two explanations were
available and they differ in consequence. Both the marginal likelihood and the squared-error
loss are bulk-dominated, so the tail --- little mass, high leverage on shape --- is weakly
constrained by either; that is a capacity limit and is fixable. Separately, the channel
destroys innovation-shape information $142$-fold against $26$-fold for the correlation
(\S\ref{ch:estimation}), so the data might simply not carry it; that is not fixable.

Varying $C$ separates them:

\begin{center}\small
\begin{tabular}{@{}rcccc@{}}
\toprule
$C$ & exact & EM--BP & local CNN & global MLP \\
\midrule
2  & $1.921$ & $-0.034$ & $1.273$ & $0.170$ \\
4  & $1.921$ & $0.812$  & $1.273$ & $0.170$ \\
8  & $1.921$ & $\mathbf{1.319}$ & $1.273$ & $0.170$ \\
12 & $1.921$ & $\mathbf{1.363}$ & $1.273$ & $0.170$ \\
16 & $1.921$ & $\mathbf{1.487}$ & $1.273$ & $0.170$ \\
\bottomrule
\end{tabular}
\end{center}

\noindent
(Generated innovation excess kurtosis, $\nseq=2048$, Laplace chain, median over three seeds,
target $1.910$. The $C=4$ entry reads $0.812$ here against $0.897$ in the four-replicate run
quoted earlier; both are the same configuration and the gap is about two standard errors of
replicate noise.) The estimator climbs monotonically, is still climbing at $C=16$, and
overtakes the network between $C=4$ and $C=8$. So the deficit is \textbf{capacity}. The
network columns are constant down the table because they do not depend on $C$ --- an internal
consistency check rather than a finding.

\begin{intuition}
The reading that survives is methodological rather than about which method wins. An estimator
can be an order of magnitude better at the task it was fitted for and still generate worse
samples, because the objective it optimises and the statistic the generated distribution
exposes weight the distribution differently. A structured method therefore has to be judged on
what it generates, which is exactly the objection that prompted this experiment.
\end{intuition}

\section{Still in progress}
\label{sec:pending}

\textsc{[pending]} Nothing above relies on this.

\begin{itemize}
\item \textbf{Does the richer kernel keep its pointwise advantage?} The component sweep varied
$C$ for generation only. If pointwise error degrades as $C$ grows, the two axes trade against
each other and the reading changes; if it does not, the estimator dominates on both. Running.
\end{itemize}

\section{Known deficiencies in the current measurements}

Recorded rather than deferred, because each affects how a number should be read.

\begin{itemize}
\item \textbf{Error bars: corrected.} An earlier version of this document stated that seeds
varied only the initialisation. That was wrong, and checking rather than assuming is what
caught it: the replicate index is mixed into every seed that draws training data, so each
replicate uses a fresh training set as well as a fresh initialisation. The held-out test set
and its exact reference are deliberately common to all replicates. The intervals therefore
cover both sampling and optimisation variance, which is the protocol one wants.
\item \textbf{One hyperparameter was selected on the evaluation set.} The convolutional
network's receptive field in \codefile{exp\_12} was chosen by an oracle over the same held-out
set the number is reported on. The bias favours the \emph{baseline}, so the quoted margin is
conservative, but the protocol is improper and the run should be repeated with a validation
split.
\item \textbf{Four quantities are asserted only in tests.} The importance-sampling validation,
the cross-family closure check, the boundary-mass diagnostic and the baseline-competence figure
are checked by the test suite but never written to an output file, so they cannot be
independently recomputed from committed data. They are marked as such wherever they appear.
\item \textbf{The generative comparison covers one family.} Four more are running.
\end{itemize}

\section{Not started}

\begin{itemize}
\item A hybrid estimator combining a learned chain kernel with a learned low-rank correction,
which is the natural next step for relaxing the Markov assumption.
\item Distillation of the fitted inference-based denoiser into a network, which is the obvious
remedy for its two-to-three-hundred-fold inference cost.
\item Anything about non-Markov priors. The misspecification limitation of
\S\ref{ch:markov} is stated and not addressed.
\end{itemize}
